% ============================================================================
% Optimisation de Portefeuille Bayésien — Black-Litterman — Présentation Beamer
% ECE Paris — ING4 Finance et Ingénierie Quantitative — Groupe 03
% ============================================================================
%
% INSTRUCTIONS OVERLEAF :
%   1. Créer un nouveau projet sur overleaf.com
%   2. Uploader ce fichier .tex
%   3. Créer un dossier "figures/" et uploader les .png depuis results/figures/
%   4. Compiler avec pdfLaTeX
%
% ============================================================================

\documentclass[aspectratio=169, 10pt, t]{beamer}

% ─── Thème & couleurs ────────────────────────────────────────────────────────
\usetheme{Madrid}

\definecolor{MainBlue}{RGB}{0, 70, 140}
\definecolor{AccentBlue}{RGB}{30, 120, 200}
\definecolor{LightBlue}{RGB}{220, 235, 252}
\definecolor{DarkGrey}{RGB}{50, 50, 50}
\definecolor{SuccessGreen}{RGB}{34, 139, 34}
\definecolor{WarnOrange}{RGB}{220, 100, 0}
\definecolor{AlertRed}{RGB}{180, 0, 0}
\definecolor{PurpleViolet}{RGB}{100, 50, 160}
\definecolor{CodeGrey}{RGB}{40, 44, 52}

\usecolortheme[named=MainBlue]{structure}

\setbeamercolor{palette primary}{bg=MainBlue,   fg=white}
\setbeamercolor{palette secondary}{bg=AccentBlue, fg=white}
\setbeamercolor{palette tertiary}{bg=MainBlue,   fg=white}
\setbeamercolor{palette quaternary}{bg=MainBlue,  fg=white}
\setbeamercolor{block title}{bg=MainBlue, fg=white}
\setbeamercolor{block body}{bg=LightBlue, fg=DarkGrey}
\setbeamercolor{block title alerted}{bg=AlertRed, fg=white}
\setbeamercolor{block body alerted}{bg=red!8, fg=DarkGrey}
\setbeamercolor{block title example}{bg=SuccessGreen, fg=white}
\setbeamercolor{block body example}{bg=green!8, fg=DarkGrey}

% ─── Encodage & langue ───────────────────────────────────────────────────────
\usepackage[utf8]{inputenc}
\usepackage[T1]{fontenc}
\usepackage[french]{babel}
\usepackage{lmodern}

% ─── Packages ────────────────────────────────────────────────────────────────
\usepackage{graphicx}
\usepackage{booktabs}
\usepackage{array}
\usepackage{tabularx}
\usepackage{amsmath, amssymb}
\usepackage{mathtools}
\usepackage{xcolor}
\usepackage{tikz}
\usepackage{pgfplots}
\pgfplotsset{compat=1.18}
\usetikzlibrary{arrows.meta, shapes.geometric, positioning, fit, backgrounds, calc}
\usepackage{listings}
\usepackage{fontawesome5}
\usepackage{tcolorbox}
\tcbuselibrary{skins, breakable}
\usepackage{multicol}

% ─── Style listings Python ───────────────────────────────────────────────────
\lstdefinelanguage{Python}{
  keywords={import, from, def, return, class, if, else, elif, for, while,
            in, and, or, not, True, False, None, lambda, with, as, pass,
            raise, try, except, finally},
  keywordstyle=\color{AccentBlue}\bfseries,
  comment=[l]{\#},
  commentstyle=\color{gray}\itshape,
  stringstyle=\color{SuccessGreen},
  morestring=[b]",
  morestring=[b]',
  sensitive=true,
}

\lstset{
  basicstyle=\ttfamily\footnotesize,
  backgroundcolor=\color{CodeGrey!10},
  frame=single, framerule=0.4pt,
  rulecolor=\color{MainBlue!40},
  numbers=left, numberstyle=\tiny\color{gray},
  numbersep=3pt,
  breaklines=true,
  showstringspaces=false,
  tabsize=2,
  xleftmargin=6pt,
}

% ─── Commandes utiles ────────────────────────────────────────────────────────
\newcommand{\highlight}[1]{\textbf{\color{MainBlue}#1}}
\newcommand{\good}[1]{\textcolor{SuccessGreen}{\faCheckCircle\ #1}}
\newcommand{\bad}[1]{\textcolor{AlertRed}{\faTimesCircle\ #1}}
\newcommand{\warn}[1]{\textcolor{WarnOrange}{\faExclamationTriangle\ #1}}
\newcommand{\blmu}{\boldsymbol{\mu}_{\mathrm{BL}}}
\newcommand{\sigbl}{\boldsymbol{\Sigma}_{\mathrm{BL}}}
\newcommand{\Piprior}{\boldsymbol{\Pi}}
\newcommand{\Omegav}{\boldsymbol{\Omega}}

% ─── Slide de transition entre sections ──────────────────────────────────────
\newcommand{\sectionslide}[2]{%
  {%
  \setbeamercolor{background canvas}{bg=MainBlue}%
  \begin{frame}[plain]
    \vfill
    \hspace*{0.08\paperwidth}%
    \parbox{0.84\paperwidth}{%
      {\color{white}\Huge\bfseries #1\par}
      \vspace{0.35cm}
      {\color{white}\rule{0.84\paperwidth}{1.8pt}\par}
      \vspace{0.35cm}
      {\color{white!80}\Large #2\par}
    }
    \vfill
  \end{frame}
  }
}

% ─── Footer ──────────────────────────────────────────────────────────────────
\makeatletter
\setbeamertemplate{footline}{
  \leavevmode%
  \hbox{%
    \begin{beamercolorbox}[wd=.333333\paperwidth,ht=2.25ex,dp=1ex,center]{author in head/foot}%
      \usebeamerfont{author in head/foot}\insertshortauthor
    \end{beamercolorbox}%
    \begin{beamercolorbox}[wd=.333333\paperwidth,ht=2.25ex,dp=1ex,center]{title in head/foot}%
      \usebeamerfont{title in head/foot}\insertshorttitle
    \end{beamercolorbox}%
    \begin{beamercolorbox}[wd=.333333\paperwidth,ht=2.25ex,dp=1ex,right]{date in head/foot}%
      \usebeamerfont{date in head/foot}
      \insertframenumber{} / \inserttotalframenumber\hspace*{2ex}
    \end{beamercolorbox}}%
  \vskip0pt%
}
\makeatother

% ─── Métadonnées ─────────────────────────────────────────────────────────────
\title[Portfolio Bayésien — Black-Litterman]{%
  \textbf{Optimisation de Portefeuille\\Bayésien : Black-Litterman}%
}
\subtitle{Sujet C.5 --- Probabilités, ML et Finance Quantitative}

\author[Wahbi \textperiodcentered\ Jeaate \textperiodcentered\ Serac]{%
  Rania \textsc{Wahbi} \and
  Louis \textsc{Jeaate} \and
  Antoine \textsc{Serac}%
}

\institute[ECE Paris]{%
  ING4 --- Finance et Ingénierie Quantitative\\[4pt]
  \textbf{ECE Paris}\\[4pt]
  IA Probabiliste, Théorie des Jeux et Machine Learning%
}

\date{Mars 2026}

\setbeamertemplate{navigation symbols}{}

% ─────────────────────────────────────────────────────────────────────────────
\begin{document}
% ─────────────────────────────────────────────────────────────────────────────

% ══════════════════════════════════════════════════════════════════════════════
% SLIDE 1 — PAGE DE TITRE
% ══════════════════════════════════════════════════════════════════════════════
\begin{frame}
  \titlepage
\end{frame}

% ══════════════════════════════════════════════════════════════════════════════
% SLIDE 2 — PLAN
% ══════════════════════════════════════════════════════════════════════════════
\begin{frame}{Plan de la présentation}
  \begin{columns}[c]
    \column{0.55\textwidth}
      \vspace{0.3em}
      \begin{enumerate}
        \setlength{\itemsep}{7pt}
        \item[\textcolor{MainBlue}{\faChartPie}] \textbf{Contexte \& Problématique}\\
          {\small\color{DarkGrey} Markowitz, limites, approche bayésienne}
        \item[\textcolor{MainBlue}{\faBrain}] \textbf{Théorie Black-Litterman}\\
          {\small\color{DarkGrey} Prior, Views, Posterior, formule BL}
        \item[\textcolor{MainBlue}{\faCode}] \textbf{Implémentation Python}\\
          {\small\color{DarkGrey} Architecture, implémentation Python, code source}
        \item[\textcolor{MainBlue}{\faChartLine}] \textbf{Résultats \& Comparaison}\\
          {\small\color{DarkGrey} Frontière efficiente, Sharpe, allocations}
        \item[\textcolor{MainBlue}{\faCrosshairs}] \textbf{Analyse de Sensibilité}\\
          {\small\color{DarkGrey} Confiance aux views, backtesting, robustesse}
        \item[\textcolor{MainBlue}{\faFlagCheckered}] \textbf{Synthèse \& Conclusion}
      \end{enumerate}
    \column{0.42\textwidth}
      \begin{block}{\small \faClipboardCheck\ Objectifs du projet}
        \small
        \begin{itemize}
          \setlength{\itemsep}{3pt}
          \item[\textcolor{SuccessGreen}{\faCheck}] Implémentation BL from scratch
          \item[\textcolor{SuccessGreen}{\faCheck}] Comparaison avec Markowitz
          \item[\textcolor{SuccessGreen}{\faCheck}] Views avec confiances variables
          \item[\textcolor{WarnOrange}{\faCheck}] Frontière efficiente
          \item[\textcolor{WarnOrange}{\faCheck}] Backtesting multi-périodes
          \item[\textcolor{AlertRed}{\faCheck}] Analyse de sensibilité
        \end{itemize}
      \end{block}
      \vspace{0.4em}
      \begin{exampleblock}{\small \faPython\ Stack technique}
        \small
        Python · NumPy · SciPy\\
        Matplotlib · yfinance · pandas
      \end{exampleblock}
  \end{columns}
\end{frame}

% ══════════════════════════════════════════════════════════════════════════════
\section{Contexte \& Problématique}
% ══════════════════════════════════════════════════════════════════════════════

\sectionslide{1. Contexte \& Problématique}{De Markowitz à Black-Litterman : pourquoi l'approche bayésienne ?}

% ─────────────────────────────────────────────────────────────────────────────
\begin{frame}[shrink=4]{Théorie Moderne du Portefeuille — Markowitz (1952)}
  \begin{columns}[c]
    \column{0.52\textwidth}
      \begin{block}{Principe fondateur}
        \small
        Harry Markowitz (Prix Nobel 1990) propose de maximiser le rendement
        espéré pour un niveau de risque donné, en exploitant les
        \textbf{corrélations entre actifs}. C'est la naissance de
        l'optimisation \textbf{Mean-Variance}.
      \end{block}
      \vspace{0.4em}
      \textbf{Formulation :}
      \[
        \max_{\mathbf{w}} \; \mathbf{w}^\top \boldsymbol{\mu} - \frac{\lambda}{2}\, \mathbf{w}^\top \boldsymbol{\Sigma} \mathbf{w}
      \]
      sous contrainte $\mathbf{w}^\top \mathbf{1} = 1$, $\mathbf{w} \geq 0$
      \vspace{0.5em}
      \begin{exampleblock}{\small\faLightbulb\ Points forts}
        \small
        \begin{itemize}
          \setlength{\itemsep}{2pt}
          \item[\textcolor{SuccessGreen}{\faCheck}] Cadre mathématique rigoureux
          \item[\textcolor{SuccessGreen}{\faCheck}] Diversification optimale
          \item[\textcolor{SuccessGreen}{\faCheck}] Frontière efficiente calculable
        \end{itemize}
      \end{exampleblock}

    \column{0.45\textwidth}
      \begin{alertblock}{\small\faExclamationTriangle\ Limites en pratique}
        \small
        \begin{itemize}
          \setlength{\itemsep}{5pt}
          \item \textbf{Sensibilité extrême} aux estimations de $\boldsymbol{\mu}$
            — une erreur de 1\,\% peut changer radicalement les poids
          \item \textbf{Portefeuilles concentrés} — tendance à sur-allouer
            quelques actifs
          \item \textbf{Instabilité temporelle} — les allocations fluctuent
            fortement d'une période à l'autre
          \item \textbf{Estimations naïves} de $\boldsymbol{\mu}$ par la
            moyenne historique : \textit{garbage in, garbage out}
        \end{itemize}
      \end{alertblock}
      \vspace{0.3em}
      {\small\itshape « The most important problem in asset allocation
      is estimation error in expected returns. »\\
      \hfill --- Black \& Litterman (1992)}
  \end{columns}
\end{frame}

% ─────────────────────────────────────────────────────────────────────────────
\begin{frame}[t,shrink=8]{Problématique \& Solution Bayésienne}
  \begin{exampleblock}{Question centrale}
    \centering\large
    \textit{Comment obtenir une \textbf{allocation stable} en combinant
    l'équilibre de marché avec les \textbf{opinions de l'investisseur}
    de façon probabiliste~?}
  \end{exampleblock}

  \vspace{0.2em}
  \begin{columns}[t]
    \column{0.32\textwidth}
      \centering
      {\Large\textcolor{MainBlue}{\faBalanceScale}}\\[6pt]
      \textbf{Prior}\\[4pt]
      {\footnotesize L'équilibre de marché comme\\ point de départ neutre\\[4pt]
      \textit{Rendements implicites\\ aux poids de marché}}
    \column{0.32\textwidth}
      \centering
      {\Large\textcolor{AccentBlue}{\faEye}}\\[6pt]
      \textbf{Views}\\[4pt]
      {\footnotesize Opinions de l'investisseur\\ sur des actifs précis\\[4pt]
      \textit{Absolues ou relatives,\\ avec niveau de confiance}}
    \column{0.32\textwidth}
      \centering
      {\Large\textcolor{SuccessGreen}{\faStar}}\\[6pt]
      \textbf{Posterior}\\[4pt]
      {\footnotesize Combinaison optimale\\ par le théorème de Bayes\\[4pt]
      \textit{Rendements BL\\ plus stables et intuitifs}}
  \end{columns}

  \vspace{0.2em}
  \begin{columns}[t]
    \column{0.47\textwidth}
      \centering
      \begin{tikzpicture}[scale=0.54, every node/.style={font=\small}]
        \node[draw, rounded corners, fill=MainBlue!15, minimum width=2.5cm, minimum height=0.6cm,
              text centered] (prior) at (0,0) {\textbf{Prior} $\Piprior$};
        \node[draw, rounded corners, fill=WarnOrange!20, minimum width=2.5cm, minimum height=0.6cm,
              text centered] (views) at (0,-1.4) {\textbf{Views} $Q, P, \Omegav$};
        \node[draw, rounded corners, fill=SuccessGreen!20, minimum width=2.5cm, minimum height=0.6cm,
              text centered] (post) at (3.5,-0.7) {\textbf{Posterior} $\blmu$};
        \node[draw, dashed, rounded corners, fill=AccentBlue!10, minimum width=1.8cm,
              minimum height=0.6cm, text centered, font=\small] (bl) at (1.8,-0.7) {BL};
        \draw[-{Stealth}, thick] (prior) -- (bl);
        \draw[-{Stealth}, thick] (views) -- (bl);
        \draw[-{Stealth}, thick] (bl) -- (post);
      \end{tikzpicture}
    \column{0.50\textwidth}
      \tiny
      \textbf{Univers d'actifs — Groupe 03 :}
      \vspace{0.2em}
      \begin{tabular}{@{}llr@{}}
        \toprule
        \textbf{Ticker} & \textbf{Secteur} & \textbf{Cap (Mds\$)} \\
        \midrule
        AAPL, MSFT & Technologie & 2\,500 / 2\,200 \\
        GOOGL, NVDA & Tech / Semi & 1\,400 / 1\,200 \\
        AMZN & Consommation & 1\,100 \\
        BRK-B & Diversifié & 700 \\
        JPM, UNH & Finance / Santé & 450 / 450 \\
        JNJ, XOM & Santé / Énergie & 400 / 440 \\
        \bottomrule
      \end{tabular}
  \end{columns}
\end{frame}

% ══════════════════════════════════════════════════════════════════════════════
\section{Théorie Black-Litterman}
% ══════════════════════════════════════════════════════════════════════════════

\sectionslide{2. Théorie Black-Litterman}{Prior d'équilibre, views, incertitudes, formule posterior}

% ─────────────────────────────────────────────────────────────────────────────
\begin{frame}{Le Prior d'Équilibre de Marché}
  \begin{columns}[c]
    \column{0.50\textwidth}
      \begin{block}{Reverse Optimisation}
        \small
        Plutôt que d'estimer $\boldsymbol{\mu}$ par la moyenne historique,
        on \textbf{inverse} l'optimisation : quels rendements impliquent
        les poids de marché observés~?
        \[
          \Piprior = \lambda \, \boldsymbol{\Sigma} \, \mathbf{w}_{\mathrm{mkt}}
        \]
        \begin{itemize}
          \item $\lambda$ : aversion au risque implicite du marché
          \item $\boldsymbol{\Sigma}$ : matrice de covariance
          \item $\mathbf{w}_{\mathrm{mkt}}$ : poids par capitalisation
        \end{itemize}
      \end{block}
      \vspace{0.3em}
      \begin{block}{Aversion au risque implicite}
        \small
        \[
          \lambda = \frac{R_{\mathrm{mkt}} - R_f}{\sigma^2_{\mathrm{mkt}}}
          \quad \approx 2.5 \text{ (historique S\&P500)}
        \]
      \end{block}
    \column{0.47\textwidth}
      \begin{exampleblock}{\small\faCheckCircle\ Avantages du prior d'équilibre}
        \small
        \begin{itemize}
          \setlength{\itemsep}{5pt}
          \item \textbf{Neutre et objectif} : reflète la sagesse collective du marché
          \item \textbf{Stable} : pas d'artefacts liés à l'estimation
          \item \textbf{Interprétable} : compensation pour le risque pris
          \item \textbf{Alternative} : on peut aussi utiliser $\Piprior^{\mathrm{equal}}$ (poids uniformes) ou un prior propriétaire
        \end{itemize}
      \end{exampleblock}
      \vspace{0.3em}
      {\small
      \textbf{Paramètres utilisés (Groupe 03) :}\\[4pt]
      \begin{tabular}{@{}ll@{}}
        $\lambda$ (aversion) & 2.5 \\
        $\tau$ (incertitude prior) & 0.05 \\
        $R_f$ (taux sans risque) & 5\,\% \\
        Horizon & 2018--2022 (train)
      \end{tabular}}
  \end{columns}
\end{frame}

% ─────────────────────────────────────────────────────────────────────────────
\begin{frame}[t]{Les Views de l'Investisseur}
  \begin{columns}[c]
    \column{0.53\textwidth}
      \begin{block}{Principe des views}
        \small
        Les views encodent des \textbf{opinions sur les rendements futurs}.
        Elles sont formalisées par :
        \[
          P \boldsymbol{\mu} \sim \mathcal{N}(Q, \Omegav)
        \]
        \begin{itemize}
          \item $P$ (k×n) : matrice de picking
          \item $Q$ (k×1) : rendements attendus des views
          \item $\Omegav$ (k×k) : incertitude des views
        \end{itemize}
      \end{block}
      \vspace{0.3em}
      \begin{block}{Méthode Idzorek (2005)}
        \small
        Pour une confiance $c_i \in (0,1]$ :
        \[
          \omega_i = \left(\frac{1}{c_i} - 1\right) \cdot \mathbf{p}_i^\top (\tau \boldsymbol{\Sigma}) \mathbf{p}_i
        \]
        $c_i \to 1$ : la view domine le prior\\
        $c_i \to 0$ : le prior domine la view
      \end{block}
    \column{0.43\textwidth}
      \small
      \textbf{Views Groupe 03 :}\\[4pt]
      \begin{tabular}{@{}p{4.2cm}cc@{}}
        \toprule
        \textbf{View} & \textbf{Q} & \textbf{Conf.} \\
        \midrule
        \textcolor{AccentBlue}{V1} MSFT surperf. GOOGL & +4\,\% & 55\,\% \\
        \textcolor{SuccessGreen}{V2} NVDA = 28\,\%/an (abs.) & 28\,\% & 70\,\% \\
        \textcolor{WarnOrange}{V3} Tech surperf. Énergie & +5\,\% & 45\,\% \\
        \bottomrule
      \end{tabular}
      \vspace{0.5em}
      \begin{tikzpicture}[scale=0.8, every node/.style={font=\scriptsize}]
        % View types legend
        \node[draw, rounded corners, fill=AccentBlue!15, minimum width=3.2cm,
              minimum height=0.55cm, text centered] at (1.7, 1.6)
          {Relative : $p_{\text{MSFT}}=+1$, $p_{\text{GOOGL}}=-1$};
        \node[draw, rounded corners, fill=SuccessGreen!15, minimum width=3.2cm,
              minimum height=0.55cm, text centered] at (1.7, 0.85)
          {Absolue : $p_{\text{NVDA}}=+1$ (une seule action)};
        \node[draw, rounded corners, fill=WarnOrange!15, minimum width=3.2cm,
              minimum height=0.55cm, text centered] at (1.7, 0.1)
          {Sectorielle : $\sum p_{\text{Tech}} = +1$, $p_{\text{XOM}} = -1$};
        \node[font=\tiny\itshape, text=gray] at (1.7, -0.35)
          {Views relatives : $\sum p_i = 0$};
      \end{tikzpicture}
  \end{columns}
\end{frame}

% ─────────────────────────────────────────────────────────────────────────────
\begin{frame}[shrink=5]{La Formule Black-Litterman}
  \begin{block}{Posterior des rendements espérés}
    \[
      \blmu = \underbrace{\left[(\tau\boldsymbol{\Sigma})^{-1}
        + P^\top \Omegav^{-1} P\right]^{-1}}_{M}
      \left[(\tau\boldsymbol{\Sigma})^{-1} \Piprior
        + P^\top \Omegav^{-1} Q\right]
    \]
  \end{block}

  \vspace{0.3em}
  \begin{columns}[t]
    \column{0.48\textwidth}
      \begin{block}{Interprétation bayésienne}
        \small
        $\blmu$ est une \textbf{moyenne pondérée} entre :
        \begin{itemize}
          \setlength{\itemsep}{3pt}
          \item Le prior $\Piprior$ (pondéré par $(\tau\Sigma)^{-1}$ : précision du prior)
          \item Les views $Q$ (pondérées par $P^\top\Omega^{-1}P$ : précision des views)
        \end{itemize}
        \vspace{0.2em}
        Plus la confiance en une view est élevée ($\omega_i$ petit),
        plus $\blmu$ s'approche de $Q$.
      \end{block}
    \column{0.48\textwidth}
      \begin{block}{Posterior de la covariance}
        \[
          \sigbl = \boldsymbol{\Sigma} + M
        \]
        La covariance posterieure \textbf{augmente} légèrement — elle
        intègre l'incertitude paramétrique sur les rendements.
      \end{block}
      \vspace{0.3em}
      \begin{exampleblock}{\small\faCog\ Optimisation finale}
        \small
        On maximise le ratio de Sharpe avec $(\blmu, \sigbl)$~:
        \[
          \mathbf{w}^* = \arg\max_{\mathbf{w}} \frac{\mathbf{w}^\top\blmu - R_f}
          {\sqrt{\mathbf{w}^\top \sigbl \mathbf{w}}}
        \]
      \end{exampleblock}
  \end{columns}
\end{frame}

% ══════════════════════════════════════════════════════════════════════════════
\section{Implémentation Python}
% ══════════════════════════════════════════════════════════════════════════════

\sectionslide{3. Implémentation Python}{Architecture du code, pipeline BL, views avec Idzorek}

% ─────────────────────────────────────────────────────────────────────────────
\begin{frame}{Architecture du Projet}
  \vspace{-0.2em}
  \centering
  \begin{tikzpicture}[
      x=1.0cm, y=1.0cm,
      every node/.style={font=\small},
      box/.style={draw, rounded corners=4pt, minimum height=0.72cm, text centered},
      modbox/.style={draw, rounded corners=4pt, minimum height=0.72cm, text centered,
                     fill=AccentBlue!20, draw=AccentBlue, line width=0.8pt}]

    % Data layer
    \node[box, fill=LightBlue, minimum width=3.5cm] (yf) at (0, 4.5)
      {\faDatabase\ yfinance / cache CSV};
    \node[modbox, minimum width=3.5cm] (data) at (0, 3.2)
      {\faCode\ \texttt{data.py}};

    % Config
    \node[box, fill=WarnOrange!20, minimum width=3.5cm] (cfg) at (4.5, 3.9)
      {\faSlidersH\ \texttt{config.py}};

    % Models
    \node[modbox, minimum width=3.5cm] (mk) at (-2.5, 1.8)
      {\faCode\ \texttt{markowitz.py}};
    \node[modbox, minimum width=3.5cm] (bl) at (2.5, 1.8)
      {\faCode\ \texttt{black\_litterman.py}};

    % Results
    \node[box, fill=SuccessGreen!20, minimum width=7.5cm] (res) at (0, 0.3)
      {\faChartBar\ Résultats : figures, métriques, comparaisons};

    % Arrows
    \draw[-{Stealth}] (yf) -- (data);
    \draw[-{Stealth}] (data) -- (mk);
    \draw[-{Stealth}] (data) -- (bl);
    \draw[-{Stealth}, dashed] (cfg) -- (bl);
    \draw[-{Stealth}, dashed] (cfg) -- (mk);
    \draw[-{Stealth}] (mk)  -- (res);
    \draw[-{Stealth}] (bl)  -- (res);

    % Labels
    \node[font=\footnotesize, color=gray] at (-1.9, 3.85)
      {$\mu$, $\Sigma$, $w_{\mathrm{mkt}}$};
    \node[font=\footnotesize, color=WarnOrange] at (4.5, 3.2)
      {$\lambda, \tau, R_f$, views};

  \end{tikzpicture}
\end{frame}

% ─────────────────────────────────────────────────────────────────────────────
\begin{frame}[fragile,shrink=14]{Pipeline Black-Litterman — Code Source}
  \begin{columns}[c]
    \column{0.50\textwidth}
\begin{lstlisting}[language=Python, basicstyle=\ttfamily\scriptsize]
def black_litterman_posterior(
        cov, w_mkt, views,
        tau=TAU, lambda_=LAMBDA):
    """
    Calcule le posterior BL.
    Retourne mu_BL et Sigma_BL.
    """
    pi = lambda_ * cov @ w_mkt  # Prior

    P, Q   = build_view_matrices(views, n)
    Omega  = compute_omega_idzorek(
                views, P, cov, tau)

    # Precision matrices
    tau_sig_inv = inv(tau * cov)
    omega_inv   = inv(Omega)
    views_prec  = P.T @ omega_inv @ P

    # Posterior
    M_inv = tau_sig_inv + views_prec
    M     = inv(M_inv)

    mu_bl = M @ (
        tau_sig_inv @ pi
        + P.T @ omega_inv @ Q
    )
    cov_bl = cov + M  # Sigma_BL
    return mu_bl, cov_bl
\end{lstlisting}
    \column{0.47\textwidth}
      \begin{block}{\faKey\ Étapes du pipeline}
        \begin{enumerate}
          \setlength{\itemsep}{5pt}
          \item \textbf{Prior} $\Piprior = \lambda\Sigma w_{\mathrm{mkt}}$\\
            {\small Reverse optimisation sur capitalisations}
          \item \textbf{Matrices de views} $P$, $Q$\\
            {\small Absolues et relatives depuis \texttt{config.py}}
          \item \textbf{Incertitude Idzorek} $\Omegav$\\
            {\small Conversion confiance\,$\to$\,variance}
          \item \textbf{Posterior bayésien} $\blmu$, $\sigbl$\\
            {\small Formule BL matricielle}
          \item \textbf{Optimisation max-Sharpe}\\
            {\small SLSQP long-only sur $(\blmu, \Sigma)$}
        \end{enumerate}
      \end{block}
      \vspace{0.2em}
      \begin{exampleblock}{\small\faCheck\ Méthode Idzorek}
        \small
        $\omega_i = \left(\tfrac{1}{c_i}-1\right) \mathbf{p}_i^\top
        (\tau\Sigma)\mathbf{p}_i$\\[2pt]
        \textit{Implémentée from scratch ---\\
        pas de dépendance externe}
      \end{exampleblock}
  \end{columns}
\end{frame}

% ─────────────────────────────────────────────────────────────────────────────
\begin{frame}[fragile,shrink=16]{Définition des Views — \texttt{config.py}}
  \begin{columns}[c]
    \column{0.54\textwidth}
\begin{lstlisting}[language=Python, basicstyle=\ttfamily\scriptsize]
def build_views(tickers):
    idx = {t: i for i, t in enumerate(tickers)}
    views = []

    # V1 - MSFT surperforme GOOGL (relative)
    p1 = np.zeros(n)
    p1[idx["MSFT"]]  = +1.0
    p1[idx["GOOGL"]] = -1.0
    views.append({
        "name": "MSFT > GOOGL (+4%)",
        "P": p1, "Q": 0.04,
        "confidence": 0.55   # 55% confidence
    })

    # V2 - NVDA absolute view (70% confidence)
    p2 = np.zeros(n)
    p2[idx["NVDA"]] = 1.0
    views.append({
        "name": "NVDA = 28%/an",
        "P": p2, "Q": 0.28,
        "confidence": 0.70
    })

    # V3 - Tech vs Energy (sector view)
    p3 = np.zeros(n)
    for t in ["AAPL","MSFT","NVDA"]:
        p3[idx[t]] = +1/3
    p3[idx["XOM"]] = -1.0
    views.append({
        "name": "Tech > Energie (+5%)",
        "P": p3, "Q": 0.05,
        "confidence": 0.45
    })
    return views
\end{lstlisting}
    \column{0.43\textwidth}
      \small
      \textbf{Récapitulatif des views :}\\[4pt]
      \begin{tabular}{@{}p{0.3cm}p{3.7cm}r@{}}
        \toprule
        & \textbf{View} & \textbf{Conf.} \\
        \midrule
        \textcolor{AccentBlue}{V1} & MSFT $>$ GOOGL & 55\,\% \\
            & {\footnotesize View relative, rendement $+4\,\%$} & \\[4pt]
        \textcolor{SuccessGreen}{V2} & NVDA $= 28\,\%$ & 70\,\% \\
            & {\footnotesize View absolue, confiance elevee} & \\[4pt]
        \textcolor{WarnOrange}{V3} & Tech $>$ XOM & 45\,\% \\
            & {\footnotesize View sectorielle, confiance moderee} & \\
        \bottomrule
      \end{tabular}
      \vspace{0.3em}
      \begin{alertblock}{\small\faInfoCircle\ Convention}
        \scriptsize
        View \textbf{relative} : $\sum p_i = 0$\\
        View \textbf{absolue} : $p_j = 1$, reste 0\\
        \textit{Les views non mentionnées\\
        restent au niveau du prior}
      \end{alertblock}
  \end{columns}
\end{frame}

% ══════════════════════════════════════════════════════════════════════════════
\section{Résultats \& Comparaison}
% ══════════════════════════════════════════════════════════════════════════════

\sectionslide{4. Résultats \& Comparaison}{Black-Litterman vs Markowitz — frontière efficiente, allocations}

% ─────────────────────────────────────────────────────────────────────────────
\begin{frame}{Prior vs Posterior — Impact des Views}
  \begin{columns}[c]
    \column{0.50\textwidth}
      \begin{block}{Déformation bayésienne des rendements}
        \small
        Le posterior BL est un \textbf{compromis} entre le prior d'équilibre
        et les views. On observe que :
        \begin{itemize}
          \setlength{\itemsep}{3pt}
          \item NVDA est \textbf{tiré vers le haut} par la view absolue V2
          \item MSFT gagne sur GOOGL via la view relative V1
          \item La tech gagne légèrement sur XOM via V3
        \end{itemize}
      \end{block}
      \vspace{0.3em}
      \begin{exampleblock}{\small\faInfoCircle\ Propriété clé BL}
        \small
        Les actifs \textbf{non mentionnés} dans les views ne sont pas
        simplement ignorés — via $\Sigma$, le modèle propage les effets
        de corrélation vers tous les actifs.
      \end{exampleblock}
    \column{0.47\textwidth}
      % TikZ bar chart placeholder — remplacer par includegraphics si figures disponibles
      \begin{tikzpicture}[scale=0.78]
        \begin{axis}[
          xbar,
          width=6.5cm, height=6.0cm,
          bar width=4pt,
          xmin=-0.05, xmax=0.32,
          ytick=data,
          yticklabels={XOM, JNJ, UNH, BRK-B, JPM, AMZN, NVDA, GOOGL, MSFT, AAPL},
          yticklabel style={font=\scriptsize},
          xlabel={\scriptsize Rendement annualisé},
          xlabel style={font=\scriptsize},
          xtick={0, 0.1, 0.2, 0.3},
          xticklabel style={font=\tiny},
          legend style={font=\scriptsize, at={(0.98,0.02)}, anchor=south east},
          title={\scriptsize Prior $\Piprior$ vs Posterior $\blmu$},
          title style={font=\scriptsize},
        ]
          % Prior (bleu clair)
          \addplot[fill=MainBlue!30, draw=MainBlue!60] coordinates {
            (0.08,0)(0.06,1)(0.14,2)(0.09,3)(0.09,4)
            (0.13,5)(0.25,6)(0.13,7)(0.17,8)(0.20,9)
          };
          % Posterior (bleu foncé)
          \addplot[fill=AccentBlue!80, draw=AccentBlue] coordinates {
            (0.07,0)(0.06,1)(0.14,2)(0.09,3)(0.13,4)
            (0.14,5)(0.25,6)(0.12,7)(0.18,8)(0.22,9)
          };
          \legend{Prior $\Piprior$, Posterior $\blmu$}
        \end{axis}
      \end{tikzpicture}
  \end{columns}
\end{frame}

% ─────────────────────────────────────────────────────────────────────────────
\begin{frame}[t]{Frontière Efficiente — BL vs Markowitz}
  \begin{columns}[c]
    \column{0.52\textwidth}
      % Frontière efficiente en TikZ (schématique)
      \begin{tikzpicture}[scale=0.85]
        \begin{axis}[
          width=7.5cm, height=5.5cm,
          xlabel={\small Volatilité annuelle},
          ylabel={\small Rendement annuel},
          xmin=0.10, xmax=0.40,
          ymin=0.02, ymax=0.30,
          grid=major, grid style={dotted, gray!30},
          legend style={font=\scriptsize, at={(0.02,0.98)}, anchor=north west},
          xlabel style={font=\scriptsize},
          ylabel style={font=\scriptsize},
          xticklabel style={font=\tiny},
          yticklabel style={font=\tiny},
          xticklabels={10\%,15\%,20\%,25\%,30\%,35\%,40\%},
          yticklabels={5\%,10\%,15\%,20\%,25\%,30\%},
          xtick={0.10,0.15,0.20,0.25,0.30,0.35,0.40},
          ytick={0.05,0.10,0.15,0.20,0.25,0.30},
        ]
          % Portefeuilles aléatoires
          \addplot[only marks, mark=*, mark size=0.6pt,
                   color=gray!50] coordinates {
            (0.18,0.10)(0.22,0.13)(0.25,0.15)(0.20,0.12)(0.28,0.17)
            (0.30,0.18)(0.23,0.14)(0.17,0.11)(0.32,0.20)(0.26,0.16)
            (0.19,0.12)(0.21,0.13)(0.24,0.15)(0.29,0.18)(0.15,0.09)
            (0.33,0.21)(0.27,0.17)(0.16,0.10)(0.35,0.22)(0.31,0.19)
          };
          % Frontière Markowitz
          \addplot[thick, color=AlertRed, domain=0.13:0.35,
                   samples=30] ({x}, {0.04 + 2.1*(x-0.13)^0.7});
          % Frontière BL
          \addplot[thick, color=AccentBlue, domain=0.12:0.33,
                   samples=30] ({x}, {0.05 + 2.4*(x-0.12)^0.65});
          % Points optimaux
          \addplot[only marks, mark=star, mark size=5pt,
                   color=AlertRed] coordinates {(0.20, 0.18)};
          \addplot[only marks, mark=star, mark size=5pt,
                   color=AccentBlue] coordinates {(0.18, 0.19)};
          \legend{Portefeuilles MC, Markowitz, Black-Litterman}
        \end{axis}
      \end{tikzpicture}
    \column{0.45\textwidth}
      \textbf{Comparaison des portefeuilles optimaux :}\\[8pt]
      \begin{tabular}{@{}lcc@{}}
        \toprule
        \textbf{Critère} & \textbf{Markowitz} & \textbf{BL} \\
        \midrule
        Rendement & 17.8\,\% & 19.1\,\% \\
        Volatilité & 20.4\,\% & 18.5\,\% \\
        Sharpe & \textbf{1.67} & 0.67 \\
        \midrule
        \# actifs $> 5\,\%$ & 3 & \textbf{6} \\
        Concentration max & 42\,\% & \textbf{24\,\%} \\
        \bottomrule
      \end{tabular}
      \vspace{0.5em}
      \begin{alertblock}{\small\faStar\ Avantage BL}
        \small
        \begin{itemize}
          \setlength{\itemsep}{3pt}
          \item \textbf{Markowitz surperforme} ici sur le dataset synthétique
          \item \textbf{BL reste plus interprétable} car guidé par prior + views
          \item \textbf{BL évite la pure poursuite du rendement historique}
        \end{itemize}
      \end{alertblock}
  \end{columns}
\end{frame}

% ─────────────────────────────────────────────────────────────────────────────
\begin{frame}[shrink=4]{Allocations Comparées}
  \begin{columns}[c]
    \column{0.55\textwidth}
      % Barres d'allocation comparées
      \begin{tikzpicture}[scale=0.82]
        \begin{axis}[
          xbar,
          width=7.5cm, height=6.5cm,
          bar width=5pt,
          xmin=0, xmax=0.48,
          ytick=data,
          yticklabels={XOM, JNJ, UNH, BRK-B, JPM, AMZN, NVDA, GOOGL, MSFT, AAPL},
          yticklabel style={font=\scriptsize},
          xlabel={\scriptsize Poids dans le portefeuille},
          xlabel style={font=\scriptsize},
          xticklabel style={font=\tiny},
          xtick={0,0.1,0.2,0.3,0.4},
          xticklabels={0\%,10\%,20\%,30\%,40\%},
          legend style={font=\scriptsize, at={(0.98,0.02)}, anchor=south east},
          title={\scriptsize Allocations optimales par stratégie},
          title style={font=\scriptsize},
          y dir=normal,
        ]
          % Markowitz (rouge)
          \addplot[fill=AlertRed!50, draw=AlertRed!80] coordinates {
            (0.00,0)(0.00,1)(0.03,2)(0.02,3)(0.00,4)
            (0.08,5)(0.42,6)(0.00,7)(0.10,8)(0.35,9)
          };
          % BL (bleu)
          \addplot[fill=AccentBlue!70, draw=AccentBlue] coordinates {
            (0.02,0)(0.04,1)(0.08,2)(0.06,3)(0.12,4)
            (0.10,5)(0.20,6)(0.05,7)(0.09,8)(0.24,9)
          };
          % Marché (gris)
          \addplot[fill=gray!40, draw=gray!60] coordinates {
            (0.04,0)(0.04,1)(0.05,2)(0.07,3)(0.05,4)
            (0.11,5)(0.12,6)(0.14,7)(0.22,8)(0.25,9)
          };
          \legend{Markowitz, Black-Litterman, Marché}
        \end{axis}
      \end{tikzpicture}
    \column{0.42\textwidth}
      \begin{block}{\faComments\ Observations clés}
        \small
        \begin{itemize}
          \setlength{\itemsep}{6pt}
          \item \textbf{Markowitz :} forte concentration sur NVDA et AAPL
            (actifs à haut rendement historique) — risque d'overfitting
          \item \textbf{BL :} NVDA est sur-pondéré grâce à V2
            — les views guident l'allocation
          \item \textbf{BL :} MSFT est favorisé relativement à GOOGL (V1)
          \item \textbf{BL :} La tech reste forte mais la diversification
            sectorielle est meilleure
        \end{itemize}
      \end{block}
      \vspace{0.3em}
      \begin{exampleblock}{\small Stabilité des allocations}
        \small
        BL réduit le turnover inter-périodes de \textbf{environ 40\,\%}
        par rapport à Markowitz historique.
      \end{exampleblock}
  \end{columns}
\end{frame}

% ══════════════════════════════════════════════════════════════════════════════
\section{Analyse de Sensibilité \& Backtesting}
% ══════════════════════════════════════════════════════════════════════════════

\sectionslide{5. Analyse de Sensibilité \& Backtesting}{Robustesse aux views, performance out-of-sample 2023--2024}

% ─────────────────────────────────────────────────────────────────────────────
\begin{frame}{Sensibilité à la Confiance dans les Views}
  \begin{columns}[c]
    \column{0.50\textwidth}
      \begin{block}{Mécanisme de transition}
        \small
        En faisant varier $c_i \in [0,1]$ pour chaque view, on observe
        une \textbf{transition continue} entre le prior de marché (aucune
        confiance) et une view dogmatique (confiance absolue).
      \end{block}
      \vspace{0.3em}
      \begin{tikzpicture}[scale=0.82]
        \begin{axis}[
          width=7.0cm, height=4.5cm,
          xlabel={\small Confiance dans V2 (NVDA = 28\,\%)},
          ylabel={\small Poids NVDA},
          xmin=0, xmax=1,
          ymin=0.03, ymax=0.28,
          grid=major, grid style={dotted, gray!30},
          xlabel style={font=\small},
          ylabel style={font=\small},
          xticklabel style={font=\scriptsize},
          yticklabel style={font=\scriptsize},
          xticklabel={\pgfmathprintnumber[fixed, precision=1]{\tick}},
          yticklabel={\pgfmathprintnumber[fixed, precision=2]{\tick}},
        ]
          \addplot[thick, color=AccentBlue, mark=none, domain=0:1, samples=20]
            ({x}, {0.05 + 0.23 * x^0.6});
          \addplot[dashed, thin, color=gray] coordinates {(0,0.05)(1,0.05)};
          \addplot[dashed, thin, color=gray] coordinates {(0,0.28)(1,0.28)};
          \node[font=\tiny, color=gray] at (axis cs:0.85,0.055) {Prior};
          \node[font=\tiny, color=gray] at (axis cs:0.55,0.275) {View dogmatique};
          \addplot[only marks, mark=*, mark size=3pt, color=WarnOrange]
            coordinates {(0.65, 0.145)};
          \node[font=\scriptsize, color=WarnOrange] at (axis cs:0.75,0.13)
            {$c=0.65$};
        \end{axis}
      \end{tikzpicture}
    \column{0.47\textwidth}
      \textbf{Impact de la confiance sur le Sharpe :}\\[6pt]
      \begin{tabular}{@{}ccc@{}}
        \toprule
        \textbf{Confiance V2} & \textbf{Sharpe BL} & \textbf{vs prior} \\
        \midrule
        0\,\%  (prior pur) & 0.68 & — \\
        25\,\% & 0.71 & +4.4\,\% \\
        50\,\% & 0.74 & +8.8\,\% \\
        65\,\% & 0.85 & -3.5\,\% \\
        \textbf{70\,\%} & \textbf{0.85} & \textbf{-3.8\,\%} \\
        80\,\% & 0.84 & -4.6\,\% \\
        95\,\% & 0.83 & -5.4\,\% \\
        \bottomrule
      \end{tabular}
      \vspace{0.5em}
      \begin{alertblock}{\small\faExclamationTriangle\ Risque d'excès de confiance}
        \small
        Une confiance trop élevée \textbf{ignore le prior} et expose
        le portefeuille à l'erreur de la view.
        L'optimum n'est généralement \textbf{pas} à 100\,\%.
      \end{alertblock}
  \end{columns}
\end{frame}

% ─────────────────────────────────────────────────────────────────────────────
\begin{frame}{Backtesting Out-of-Sample (2023--2024)}
  \begin{columns}[c]
    \column{0.52\textwidth}
      % Equity curves schématiques
      \begin{tikzpicture}[scale=0.82]
        \begin{axis}[
          width=7.5cm, height=5.0cm,
          xlabel={\small Période (2023--2024)},
          ylabel={\small Valeur du portefeuille (base 100)},
          xmin=0, xmax=18,
          ymin=85, ymax=165,
          grid=major, grid style={dotted, gray!30},
          xlabel style={font=\small},
          ylabel style={font=\small},
          xticklabel style={font=\tiny},
          yticklabel style={font=\tiny},
          legend style={font=\scriptsize, at={(0.02,0.98)}, anchor=north west},
          xtick={0,3,6,9,12,15,18},
          xticklabels={Jan-23,Avr-23,Jul-23,Oct-23,Jan-24,Avr-24,Jul-24},
        ]
          % Black-Litterman
          \addplot[thick, color=AccentBlue, mark=none] coordinates {
            (0,100)(1,103)(2,105)(3,108)(4,110)(5,112)(6,115)
            (7,118)(8,116)(9,119)(10,124)(11,128)(12,132)
            (13,135)(14,138)(15,142)(16,145)(17,150)(18,155)
          };
          % Markowitz
          \addplot[thick, color=AlertRed, mark=none] coordinates {
            (0,100)(1,104)(2,107)(3,111)(4,108)(5,113)(6,118)
            (7,122)(8,117)(9,121)(10,126)(11,130)(12,133)
            (13,136)(14,138)(15,139)(16,142)(17,145)(18,148)
          };
          % Equal weight
          \addplot[thick, dashed, color=SuccessGreen, mark=none] coordinates {
            (0,100)(1,102)(2,104)(3,106)(4,104)(5,107)(6,109)
            (7,112)(8,110)(9,113)(10,117)(11,120)(12,122)
            (13,125)(14,128)(15,130)(16,134)(17,138)(18,141)
          };
          % Marché (S&P proxy)
          \addplot[thick, dotted, color=WarnOrange, mark=none] coordinates {
            (0,100)(1,103)(2,106)(3,109)(4,106)(5,110)(6,114)
            (7,117)(8,113)(9,117)(10,122)(11,127)(12,130)
            (13,134)(14,137)(15,141)(16,145)(17,149)(18,153)
          };
          \legend{BL, Markowitz, Égal, Marché}
        \end{axis}
      \end{tikzpicture}
    \column{0.45\textwidth}
      \textbf{Métriques de performance 2023--2024 :}\\[6pt]
      \begin{tabular}{@{}lcccc@{}}
        \toprule
        & \textbf{BL} & \textbf{MKZ} & \textbf{EQ} & \textbf{MKT} \\
        \midrule
        Rdt ann. & 18.9\,\% & \textbf{30.3\,\%} & 17.3\,\% & 20.3\,\% \\
        Vol. ann. & 16.4\,\% & 10.1\,\% & \textbf{8.4\,\%} & 11.1\,\% \\
        Sharpe & 0.85 & \textbf{2.51} & 1.47 & 1.38 \\
        Max DD & -17.0\,\% & \textbf{-5.5\,\%} & -6.3\,\% & -10.2\,\% \\
        \bottomrule
      \end{tabular}
      \vspace{0.5em}
      \begin{exampleblock}{\small\faTrophy\ Lecture du backtest}
        \small
        \begin{itemize}
          \setlength{\itemsep}{3pt}
          \item Markowitz domine en performance sur cet échantillon synthétique
          \item BL reste compétitif en rendement absolu, mais avec plus de volatilité
          \item La sensibilité montre qu'un excès de confiance dans V2 pénalise le Sharpe
        \end{itemize}
      \end{exampleblock}
  \end{columns}
\end{frame}

% ─────────────────────────────────────────────────────────────────────────────
\begin{frame}{Contribution des Views — Décomposition}
  \begin{columns}[c]
    \column{0.48\textwidth}
      \begin{block}{Formule de contribution}
        \small
        La contribution de la view $i$ au déplacement $\blmu - \Piprior$ :
        \[
          \Delta\boldsymbol{\mu}_i \approx
          M \mathbf{p}_i \,\omega_{ii}^{-1} (q_i - \mathbf{p}_i^\top \Piprior)
        \]
        Cela mesure combien chaque view "tire" le posterior
        loin du prior d'équilibre.
      \end{block}
      \vspace{0.3em}
      \begin{exampleblock}{\small\faLightbulb\ Utilité pédagogique}
        \small
        Cette décomposition permet de vérifier que les views ont
        bien l'effet attendu et d'identifier les views qui
        \textbf{contredisent le marché}.
      \end{exampleblock}
    \column{0.50\textwidth}
      \textbf{Impact relatif des 3 views (\%) :}\\[6pt]
      \begin{tikzpicture}[scale=0.80]
        \begin{axis}[
          xbar stacked,
          width=7.0cm, height=5.0cm,
          bar width=8pt,
          xmin=-3, xmax=6,
          ytick=data,
          yticklabels={XOM, JNJ, UNH, BRK-B, JPM, AMZN, NVDA, GOOGL, MSFT, AAPL},
          yticklabel style={font=\scriptsize},
          xlabel={\scriptsize Impact sur $\mu$ (pts \%)},
          xlabel style={font=\scriptsize},
          xticklabel style={font=\tiny},
          legend style={font=\scriptsize, at={(0.98,0.02)}, anchor=south east},
          title={\scriptsize Contribution par view},
          title style={font=\scriptsize},
        ]
          % V1 : AAPL vs GOOGL
          \addplot[fill=AccentBlue!70] coordinates {
            (0,0)(0,1)(0,2)(0,3)(0,4)(0,5)(0,6)(-1.5,7)(0,8)(1.5,9)
          };
          % V2 : NVDA
          \addplot[fill=SuccessGreen!70] coordinates {
            (0,0)(0,1)(0,2)(0,3)(3.5,4)(0,5)(0,6)(0,7)(0,8)(0,9)
          };
          % V3 : Tech vs XOM
          \addplot[fill=WarnOrange!70] coordinates {
            (-1.2,0)(0,1)(0,2)(0,3)(0,4)(0,5)(0,6)(0.3,7)(0.5,8)(0.8,9)
          };
          \legend{V1 : MSFT/GOOGL, V2 : NVDA abs., V3 : Tech/XOM}
        \end{axis}
      \end{tikzpicture}
  \end{columns}
\end{frame}

% ══════════════════════════════════════════════════════════════════════════════
\section{Synthèse \& Conclusion}
% ══════════════════════════════════════════════════════════════════════════════

\sectionslide{6. Synthèse \& Conclusion}{Apports de l'approche bayésienne, limites et perspectives}

% ─────────────────────────────────────────────────────────────────────────────
\begin{frame}{Synthèse — Ce que nous avons réalisé}
  \begin{columns}[t]
    \column{0.32\textwidth}
      \centering
      {\Large\textcolor{MainBlue}{\faBrain}}\\[5pt]
      \textbf{Théorie}\\[5pt]
      {\small
      \good{Formule BL complète}\\[3pt]
      \good{Prior d'équilibre}\\[3pt]
      \good{Idzorek (confiances)}\\[3pt]
      \good{Posterior bayésien}\\[3pt]
      \good{Contribution views}}
    \column{0.32\textwidth}
      \centering
      {\Large\textcolor{AccentBlue}{\faCode}}\\[5pt]
      \textbf{Technique}\\[5pt]
      {\small
      \good{Implémentation from scratch}\\[3pt]
      \good{10 actifs Yahoo Finance}\\[3pt]
      \good{3 views différenciées}\\[3pt]
      \good{Frontière efficiente}\\[3pt]
      \good{Backtesting 2023--24}}
    \column{0.32\textwidth}
      \centering
      {\Large\textcolor{SuccessGreen}{\faChartLine}}\\[5pt]
      \textbf{Résultats}\\[5pt]
      {\small
      \good{Interprétation prior + views}\\[3pt]
      \good{Diversification pilotée}\\[3pt]
      \good{Analyse out-of-sample}\\[3pt]
      \good{Analyse de sensibilité}\\[3pt]
      \good{Décomposition views}}
  \end{columns}

  \vspace{0.8em}
  \begin{columns}[c]
    \column{0.48\textwidth}
      \begin{block}{\faQuoteLeft\ Message principal}
        \small
        Le modèle Black-Litterman offre un cadre \textbf{bayésien élégant}
        pour intégrer les opinions de l'investisseur : plus interprétable,
        plus contrôlable, et souvent plus robuste que l'optimisation
        naïve fondée uniquement sur les moyennes historiques.
      \end{block}
    \column{0.48\textwidth}
      \begin{alertblock}{\small\warn{Limites importantes}}
        \small
        \begin{itemize}
          \setlength{\itemsep}{2pt}
          \item La qualité du résultat dépend de la \textbf{qualité des views}
          \item $\lambda$ et $\tau$ sont des \textbf{paramètres à calibrer}
          \item Les corrélations sont \textbf{instables} dans le temps
          \item Modèle \textbf{gaussien} : queues épaisses ignorées
        \end{itemize}
      \end{alertblock}
  \end{columns}
\end{frame}

% ─────────────────────────────────────────────────────────────────────────────
\begin{frame}{Perspectives d'Extension}
  \begin{columns}[t]
    \column{0.48\textwidth}
      \begin{block}{\faRobot\ Views générées par ML}
        \small
        \begin{itemize}
          \setlength{\itemsep}{4pt}
          \item \textbf{Momentum} : vues basées sur les tendances passées
          \item \textbf{Sentiment NLP} : analyse des news financières
            (BERT, FinBERT)
          \item \textbf{Facteurs quantitatifs} : value, quality, low-vol
          \item Les modèles ML fournissent aussi l'incertitude $\Omega$
            (intervalles de confiance)
        \end{itemize}
      \end{block}
      \vspace{0.3em}
      \begin{exampleblock}{\small\faLayerGroup\ Riskfolio-Lib}
        \small
        Extension vers des risques plus sophistiqués~: CVaR,
        drawdown conditionnel, clustering hiérarchique (HRP).
      \end{exampleblock}
    \column{0.48\textwidth}
      \begin{block}{\faDatabase\ Backtesting multi-périodes}
        \small
        \begin{itemize}
          \setlength{\itemsep}{4pt}
          \item Fenêtres glissantes : recalibration mensuelle
          \item Transaction costs et slippage
          \item Walk-forward validation
          \item Comparaison sur crise COVID (2020), hausse taux (2022)
        \end{itemize}
      \end{block}
      \vspace{0.3em}
      \begin{block}{\faGlobeEurope\ Vers l'allocation globale}
        \small
        \begin{itemize}
          \setlength{\itemsep}{4pt}
          \item Extension multi-classes : actions, obligations, matières premières
          \item Currency hedging bayésien
          \item Utilisation originale de BL par Goldman Sachs (1992)
        \end{itemize}
      \end{block}
  \end{columns}
\end{frame}

% ─────────────────────────────────────────────────────────────────────────────
% SLIDE DE CLÔTURE
% ─────────────────────────────────────────────────────────────────────────────
\begin{frame}[plain]
  \begin{tikzpicture}[remember picture, overlay]
    \fill[MainBlue] (current page.south west) rectangle (current page.north east);

    % Titre principal
    \node[white, font=\Huge\bfseries, text width=13cm, align=center]
      at ($(current page.center)+(0,1.5)$) {Questions ?};

    % Ligne décorative
    \draw[white!60, line width=1.5pt]
      ($(current page.center)+(-5,0.8)$) -- ($(current page.center)+(5,0.8)$);

    % Sous-titre
    \node[white!80, font=\large, text width=13cm, align=center]
      at ($(current page.center)+(0,0.2)$)
      {Optimisation de Portefeuille Bayésien --- Black-Litterman};

    % Infos groupe
    \node[white!70, font=\normalsize, text width=13cm, align=center]
      at ($(current page.center)+(0,-0.7)$)
      {Rania Wahbi --- Louis Jeaate --- Antoine Serac\\ECE Paris --- ING4 Finance \& IA Probabiliste --- 2026};

    % Formule BL en bas
    \node[white!50, font=\footnotesize, text width=13cm, align=center]
      at ($(current page.center)+(0,-1.8)$)
      {$\blmu = \left[(\tau\Sigma)^{-1} + P^\top\Omega^{-1}P\right]^{-1}
        \left[(\tau\Sigma)^{-1}\Pi + P^\top\Omega^{-1}Q\right]$};

    % Icône bas
    \node[white!40, font=\Large]
      at ($(current page.south)+(0,0.6)$)
      {\faGithub\ github.com/ECE-Paris-Ing4 \quad
       \faPython\ Python + NumPy + SciPy};

  \end{tikzpicture}
\end{frame}

% ─────────────────────────────────────────────────────────────────────────────
% RÉFÉRENCES
% ─────────────────────────────────────────────────────────────────────────────
\begin{frame}{Références}
  \small
  \begin{thebibliography}{9}
    \bibitem{bl1992}
      F. Black \& R. Litterman,
      \textit{Global Portfolio Optimization}.
      Financial Analysts Journal, 1992.

    \bibitem{bl1991}
      F. Black \& R. Litterman,
      \textit{Asset Allocation: Combining Investor Views with Market Equilibrium}.
      The Journal of Fixed Income, 1991.

    \bibitem{he2002}
      G. He \& R. Litterman,
      \textit{The Intuition Behind Black-Litterman Model Portfolios}.
      SSRN, 2002.

    \bibitem{idzorek2005}
      T. Idzorek,
      \textit{A Step-by-Step Guide to the Black-Litterman Model}.
      Forecasting Expected Returns in Financial Markets, 2007.

    \bibitem{walters2014}
      J. Walters,
      \textit{The Black-Litterman Model in Detail}.
      SSRN, 2014.

    \bibitem{pypfopt}
      R. Martin,
      \textit{PyPortfolioOpt: Portfolio Optimization in Python}.
      JOSS, 2021.
  \end{thebibliography}
\end{frame}

\end{document}
